\chapter{Anhang}
\label{ch:anhang}

Der Anhang dokumentiert die Einstellungen der Trainingsläufe und zeigt
jene Verläufe vollständig, die \Cref{ch:ergebnisse} nur ausschnittweise
wiedergibt. Die Abbildungen decken alle 24 Matchups ab, getrennt nach
verteidigter Topologie: Crown-Jewel-Quote, gehaltene Knoten, Schritte mit
\ac{SLA}-Bruch, Abwehrbonus und Konvergenzmarken.

\section{Hyperparameter}
\label{sec:hyperparameter}

Die Agenten werden mit den Voreinstellungen von \ac{SB3} trainiert; im
Programmcode wird kein Hyperparameter gesetzt. \Cref{tab:hyperparameter} führt
die Werte auf, die daraus folgen.

Diese Zurückhaltung ist beabsichtigt. Würden die Parameter angepasst, so
geschähe das notwendigerweise an einer bestimmten Topologie, und die dort
gefundenen Werte bevorzugten diese gegenüber den übrigen. Da der Vergleich der
Topologien der Gegenstand der Arbeit ist, wäre damit genau die unabhängige
Variable verunreinigt. Einheitliche Voreinstellungen halten diesen Einfluss
konstant.

\begin{table}[H]
  \centering
  \caption[Hyperparameter des Trainings]{Hyperparameter des Trainings. Alle
  Werte entsprechen den Voreinstellungen von \ac{SB3} in der eingesetzten
  Fassung; abweichend gesetzt wird allein die Zahl der Zeitschritte.}
  \label{tab:hyperparameter}
  \begin{tabularx}{\textwidth}{llX}
    \toprule
    Parameter & Wert & Bedeutung \\
    \midrule
    \texttt{learning\_rate}       & \num{0.0003} & Schrittweite $\alpha$ in \Cref{eq:gradientenschritt} \\
    \texttt{n\_steps}             & \num{2048}   & Länge eines Rollouts, also gesammelte Schritte je Lernphase \\
    \texttt{batch\_size}          & \num{64}     & Größe eines Minibatches \\
    \texttt{n\_epochs}            & \num{10}     & Durchgänge über denselben Rollout \\
    \texttt{gamma}                & \num{0.99}   & Diskontfaktor $\gamma$ aus \Cref{eq:return} \\
    \texttt{gae\_lambda}          & \num{0.95}   & Gewichtung zwischen Critic-Schätzung und beobachteten Rewards bei der Vorteilsschätzung \\
    \texttt{clip\_range}          & \num{0.2}    & Breite des Bandes $\varepsilon$ in \Cref{eq:ppo-clip} \\
    \texttt{normalize\_advantage} & wahr         & Vorteile werden je Minibatch auf Mittelwert null und Streuung eins gebracht \\
    \texttt{ent\_coef}            & \num{0.0}    & Gewicht des Entropie-Anteils, der zusätzliche Exploration begünstigen würde \\
    \texttt{vf\_coef}             & \num{0.5}    & Gewicht des Critic-Verlusts gegenüber dem des Actors \\
    \texttt{max\_grad\_norm}      & \num{0.5}    & Obergrenze für die Länge des Gradienten je Schritt \\
    \midrule
    Zeitschritte je Lauf          & \num{500000} & eigens gesetzt, siehe \Cref{sec:phasen} \\
    \bottomrule
  \end{tabularx}
\end{table}

\section{Crown-Jewel-Quote je Topologie}
\label{app:cj}

Die folgenden vier Abbildungen zeigen den Anteil der Episoden mit erreichtem
Crown Jewel für alle 24 Matchups. \Cref{sec:cj-verlauf} stellt daraus nur die
vier Matchups auf der Diagonalen dar. Gebildet ist die Quote wie dort über ein
Fenster von elf Episoden und alle fünf Seeds; die Skala ist in allen Feldern
dieselbe. Fett überschrieben ist jeweils das Feld des Angreifers, der auf der
verteidigten Topologie trainiert wurde.

\begin{figure}[H]
  \centering
  \includegraphics[width=\textwidth]{cj_flat}
  \caption[Crown-Jewel-Quote auf \emph{flat}]{Crown-Jewel-Quote der Verteidiger
  auf \emph{flat} gegen alle sechs Angreifer.}
  \label{fig:cj-flat}
\end{figure}

\begin{figure}[H]
  \centering
  \includegraphics[width=\textwidth]{cj_hub_and_spoke}
  \caption[Crown-Jewel-Quote auf \emph{hub\_and\_spoke}]{Crown-Jewel-Quote der
  Verteidiger auf \emph{hub\_and\_spoke} gegen alle sechs Angreifer.}
  \label{fig:cj-hs}
\end{figure}

\begin{figure}[H]
  \centering
  \includegraphics[width=\textwidth]{cj_dmz}
  \caption[Crown-Jewel-Quote auf \emph{dmz}]{Crown-Jewel-Quote der Verteidiger
  auf \emph{dmz} gegen alle sechs Angreifer.}
  \label{fig:cj-dmz}
\end{figure}

\begin{figure}[H]
  \centering
  \includegraphics[width=\textwidth]{cj_micro_segmented}
  \caption[Crown-Jewel-Quote auf \emph{micro\_segmented}]{Crown-Jewel-Quote der
  Verteidiger auf \emph{micro\_segmented} gegen alle sechs Angreifer.}
  \label{fig:cj-ms}
\end{figure}

\section{Gehaltene Knoten je Topologie}
\label{app:knoten}

Die folgenden vier Abbildungen zeigen die größte Zahl gleichzeitig gehaltener
Knoten je Episode für alle 24 Matchups, gemittelt über die fünf Seeds und
geglättet über fünf Episoden. \Cref{sec:knoten-verlauf} stellt daraus nur
die vier Matchups auf der Diagonalen dar. Die Skala reicht in allen Feldern
von null bis zehn Knoten. Fett überschrieben ist jeweils das Feld des
Angreifers, der auf der verteidigten Topologie trainiert wurde.

\begin{figure}[H]
  \centering
  \includegraphics[width=\textwidth]{knoten_flat}
  \caption[Gehaltene Knoten auf \emph{flat}]{Gehaltene Knoten der Verteidiger
  auf \emph{flat} gegen alle sechs Angreifer.}
  \label{fig:knoten-flat}
\end{figure}

\begin{figure}[H]
  \centering
  \includegraphics[width=\textwidth]{knoten_hub_and_spoke}
  \caption[Gehaltene Knoten auf \emph{hub\_and\_spoke}]{Gehaltene Knoten der Verteidiger
  auf \emph{hub\_and\_spoke} gegen alle sechs Angreifer.}
  \label{fig:knoten-hs}
\end{figure}

\begin{figure}[H]
  \centering
  \includegraphics[width=\textwidth]{knoten_dmz}
  \caption[Gehaltene Knoten auf \emph{dmz}]{Gehaltene Knoten der Verteidiger
  auf \emph{dmz} gegen alle sechs Angreifer.}
  \label{fig:knoten-dmz}
\end{figure}

\begin{figure}[H]
  \centering
  \includegraphics[width=\textwidth]{knoten_micro_segmented}
  \caption[Gehaltene Knoten auf \emph{micro\_segmented}]{Gehaltene Knoten der Verteidiger
  auf \emph{micro\_segmented} gegen alle sechs Angreifer.}
  \label{fig:knoten-ms}
\end{figure}

\section{\acs{SLA}-Brüche je Topologie}
\label{app:sla}

Die folgenden vier Abbildungen zeigen, in wie vielen Schritten je Episode die
Verfügbarkeit unter der Schwelle lag, gemittelt über die fünf Seeds. Jedes Feld
trägt eine eigene Skala, weil die Werte weit auseinanderliegen; $\Sigma$ nennt
die Summe aller Bruchschritte des Matchups über alle fünf Seeds.
\Cref{fig:sla-einbrueche} zeigt zwei dieser Fälle nach Seeds getrennt.

\begin{figure}[H]
  \centering
  \includegraphics[width=\textwidth]{sla_flat}
  \caption[\acs{SLA}-Brüche auf \emph{flat}]{Schritte mit \ac{SLA}-Bruch der
  Verteidiger auf \emph{flat} gegen alle sechs Angreifer.}
  \label{fig:sla-flat}
\end{figure}

\begin{figure}[H]
  \centering
  \includegraphics[width=\textwidth]{sla_hub_and_spoke}
  \caption[\acs{SLA}-Brüche auf \emph{hub\_and\_spoke}]{Schritte mit
  \ac{SLA}-Bruch der Verteidiger auf \emph{hub\_and\_spoke} gegen alle sechs
  Angreifer.}
  \label{fig:sla-hs}
\end{figure}

\begin{figure}[H]
  \centering
  \includegraphics[width=\textwidth]{sla_dmz}
  \caption[\acs{SLA}-Brüche auf \emph{dmz}]{Schritte mit \ac{SLA}-Bruch der
  Verteidiger auf \emph{dmz} gegen alle sechs Angreifer.}
  \label{fig:sla-dmz}
\end{figure}

\begin{figure}[H]
  \centering
  \includegraphics[width=\textwidth]{sla_micro_segmented}
  \caption[\acs{SLA}-Brüche auf \emph{micro\_segmented}]{Schritte mit
  \ac{SLA}-Bruch der Verteidiger auf \emph{micro\_segmented} gegen alle sechs
  Angreifer.}
  \label{fig:sla-ms}
\end{figure}


\section{Abwehrbonus je Topologie}
\label{app:abwehr}

Die folgenden vier Abbildungen zeigen den Abwehrbonus je Episode für alle 24
Matchups, gemittelt über die fünf Seeds und geglättet über fünf Episoden;
\Cref{sec:aktionen} zeigt daraus nur das Mittel über die sechs Angreifer je
Topologie. Die Skala ist in allen Feldern dieselbe, die Nulllinie ist
eingezeichnet. Fett überschrieben ist jeweils das Feld des Angreifers, der
auf der verteidigten Topologie trainiert wurde.

\begin{figure}[H]
  \centering
  \includegraphics[width=\textwidth]{abwehr_flat}
  \caption[Abwehrbonus auf \emph{flat}]{Abwehrbonus der Verteidiger
  auf \emph{flat} gegen alle sechs Angreifer.}
  \label{fig:abwehr-flat}
\end{figure}

\begin{figure}[H]
  \centering
  \includegraphics[width=\textwidth]{abwehr_hub_and_spoke}
  \caption[Abwehrbonus auf \emph{hub\_and\_spoke}]{Abwehrbonus der Verteidiger
  auf \emph{hub\_and\_spoke} gegen alle sechs Angreifer.}
  \label{fig:abwehr-hs}
\end{figure}

\begin{figure}[H]
  \centering
  \includegraphics[width=\textwidth]{abwehr_dmz}
  \caption[Abwehrbonus auf \emph{dmz}]{Abwehrbonus der Verteidiger
  auf \emph{dmz} gegen alle sechs Angreifer.}
  \label{fig:abwehr-dmz}
\end{figure}

\begin{figure}[H]
  \centering
  \includegraphics[width=\textwidth]{abwehr_micro_segmented}
  \caption[Abwehrbonus auf \emph{micro\_segmented}]{Abwehrbonus der Verteidiger
  auf \emph{micro\_segmented} gegen alle sechs Angreifer.}
  \label{fig:abwehr-ms}
\end{figure}

\section{Konvergenzmarken je Topologie}
\label{app:konvergenz}

Die folgenden vier Abbildungen zeigen für alle 24 Matchups, an welcher Stelle
des Verlaufs die beiden in \Cref{sec:konvergenzergebnis} verglichenen Maße
greifen. Gezeichnet ist der Verteidiger-Reward, gemittelt über die fünf Seeds
und geglättet über fünf Episoden. Die drei Senkrechten markieren das Kriterium
aus \Cref{sec:konvergenz}, dasselbe Kriterium mit dem strengeren Trendtest aus
\Cref{sec:konvergenzergebnis} und die 90-Prozent-Marke aus
\Cref{sec:limitationen}, jeweils gemittelt über die fünf Seeds. Fett überschrieben ist das Feld des Angreifers, der auf
der verteidigten Topologie trainiert wurde.

\begin{figure}[H]
  \centering
  \includegraphics[width=\textwidth]{konvergenz_flat}
  \caption[Konvergenzmarken auf \emph{flat}]{Konvergenzmarken der Verteidiger
  auf \emph{flat} gegen alle sechs Angreifer.}
  \label{fig:konv-flat}
\end{figure}

\begin{figure}[H]
  \centering
  \includegraphics[width=\textwidth]{konvergenz_hub_and_spoke}
  \caption[Konvergenzmarken auf \emph{hub\_and\_spoke}]{Konvergenzmarken der
  Verteidiger auf \emph{hub\_and\_spoke} gegen alle sechs Angreifer.}
  \label{fig:konv-hs}
\end{figure}

\begin{figure}[H]
  \centering
  \includegraphics[width=\textwidth]{konvergenz_dmz}
  \caption[Konvergenzmarken auf \emph{dmz}]{Konvergenzmarken der Verteidiger
  auf \emph{dmz} gegen alle sechs Angreifer.}
  \label{fig:konv-dmz}
\end{figure}

\begin{figure}[H]
  \centering
  \includegraphics[width=\textwidth]{konvergenz_micro_segmented}
  \caption[Konvergenzmarken auf \emph{micro\_segmented}]{Konvergenzmarken der
  Verteidiger auf \emph{micro\_segmented} gegen alle sechs Angreifer.}
  \label{fig:konv-ms}
\end{figure}
